\documentclass[10pt,a4paper]{amsart}
\usepackage{amsmath,amssymb,amsthm,graphicx}
% Retained preamble for the Stacks Project chapter morphisms.tex (Morphisms of Schemes).
% The chapter file pinned-original-source.tex carries only \input{preamble} on its first line, so no
% preamble text lives inside the chapter itself.  The Stacks Project uses one single shared preamble for
% every chapter; the snapshot stages it verbatim in sources/stacks-chow/preamble.tex.  The block below is
% copied mechanically, line for line, from that file's lines 14-18 (the Xy-pic setup) and lines 154-217
% (the native theorem declarations together with the mathematical macros).  Nothing was added, removed or
% redefined.  The wrapper class and the amsmath/amssymb/amsthm/graphicx base packages are supplied by the
% assembler, so the class/\externaldocument part of the shared preamble is not repeated here.
% For commutative diagrams we use Xy-pic
\usepackage[all]{xy}

% We use 2cell for 2-commutative diagrams.
\xyoption{2cell}
% Theorem environments.
%
\theoremstyle{plain}
\newtheorem{theorem}[subsection]{Theorem}
\newtheorem{proposition}[subsection]{Proposition}
\newtheorem{lemma}[subsection]{Lemma}

\theoremstyle{definition}
\newtheorem{definition}[subsection]{Definition}
\newtheorem{example}[subsection]{Example}
\newtheorem{exercise}[subsection]{Exercise}
\newtheorem{situation}[subsection]{Situation}

\theoremstyle{remark}
\newtheorem{remark}[subsection]{Remark}
\newtheorem{remarks}[subsection]{Remarks}

\numberwithin{equation}{subsection}

% Macros
%
\def\lim{\mathop{\mathrm{lim}}\nolimits}
\def\colim{\mathop{\mathrm{colim}}\nolimits}
\def\Spec{\mathop{\mathrm{Spec}}}
\def\Hom{\mathop{\mathrm{Hom}}\nolimits}
\def\Ext{\mathop{\mathrm{Ext}}\nolimits}
\def\SheafHom{\mathop{\mathcal{H}\!\mathit{om}}\nolimits}
\def\SheafExt{\mathop{\mathcal{E}\!\mathit{xt}}\nolimits}
\def\Sch{\mathit{Sch}}
\def\Mor{\mathop{\mathrm{Mor}}\nolimits}
\def\Ob{\mathop{\mathrm{Ob}}\nolimits}
\def\Sh{\mathop{\mathit{Sh}}\nolimits}
\def\NL{\mathop{N\!L}\nolimits}
\def\CH{\mathop{\mathrm{CH}}\nolimits}
\def\proetale{{pro\text{-}\acute{e}tale}}
\def\etale{{\acute{e}tale}}
\def\QCoh{\mathit{QCoh}}
\def\Ker{\mathop{\mathrm{Ker}}}
\def\Im{\mathop{\mathrm{Im}}}
\def\Coker{\mathop{\mathrm{Coker}}}
\def\Coim{\mathop{\mathrm{Coim}}}

% Boxtimes
%
\DeclareMathSymbol{\boxtimes}{\mathbin}{AMSa}{"02}

%
% Macros for moduli stacks/spaces
%
\def\QCohstack{\mathcal{QC}\!\mathit{oh}}
\def\Cohstack{\mathcal{C}\!\mathit{oh}}
\def\Spacesstack{\mathcal{S}\!\mathit{paces}}
\def\Quotfunctor{\mathrm{Quot}}
\def\Hilbfunctor{\mathrm{Hilb}}
\def\Curvesstack{\mathcal{C}\!\mathit{urves}}
\def\Polarizedstack{\mathcal{P}\!\mathit{olarized}}
\def\Complexesstack{\mathcal{C}\!\mathit{omplexes}}
% \Pic is the operator that assigns to X its picard group, usage \Pic(X)
% \Picardstack_{X/B} denotes the Picard stack of X over B
% \Picardfunctor_{X/B} denotes the Picard functor of X over B
\def\Pic{\mathop{\mathrm{Pic}}\nolimits}
\def\Picardstack{\mathcal{P}\!\mathit{ic}}
\def\Picardfunctor{\mathrm{Pic}}
\def\Deformationcategory{\mathcal{D}\!\mathit{ef}}

% --- appended 2026-10-12 (Chief): the project own definitions of the reference/slogan/history
% display environments, copied mechanically from the shared Stacks preamble lines 8-12
% (\usepackage{verbatim} plus the three \newenvironment lines, which the chapter preamble keeps
% verbatim from the project preamble). Nothing added, removed or redefined: these environments
% carry no mathematical content and the project itself typesets them as comments. Without these
% three lines a counted statement whose own text uses them fails with Environment slogan undefined.
\usepackage{verbatim}
\newenvironment{reference}{\comment}{\endcomment}
\newenvironment{slogan}{\comment}{\endcomment}
\newenvironment{history}{\comment}{\endcomment}
\usepackage{cleveref}
\crefname{originalequation}{original equation}{original equations}
\Crefname{originalequation}{Original equation}{Original equations}
\makeatletter
\newlabel{equation-dimension-formula}{{02JV}{source}{}{sourcepointer-1}{}}
\newlabel{equation-dimension-formula@cref}{{[originalequation][0][]02JV}{[1][0][]source}{}{}{}}
\makeatother
\crefname{originallemma}{original lemma}{original lemmas}
\Crefname{originallemma}{Original lemma}{Original lemmas}
\makeatletter
\newlabel{lemma-locally-finite-type-characterize}{{01T2}{source}{}{sourcepointer-2}{}}
\newlabel{lemma-locally-finite-type-characterize@cref}{{[originallemma][0][]01T2}{[1][0][]source}{}{}{}}
\makeatother
\crefname{originallemma}{original lemma}{original lemmas}
\Crefname{originallemma}{Original lemma}{Original lemmas}
\makeatletter
\newlabel{lemma-universally-catenary-local}{{02J9}{source}{}{sourcepointer-3}{}}
\newlabel{lemma-universally-catenary-local@cref}{{[originallemma][0][]02J9}{[1][0][]source}{}{}{}}
\makeatother
\crefname{originallemma}{original lemma}{original lemmas}
\Crefname{originallemma}{Original lemma}{Original lemmas}
\makeatletter
\newlabel{properties-lemma-codimension-local-ring}{{02IZ}{source}{}{sourcepointer-4}{}}
\newlabel{properties-lemma-codimension-local-ring@cref}{{[originallemma][0][]02IZ}{[1][0][]source}{}{}{}}
\makeatother
\crefname{originallemma}{original lemma}{original lemmas}
\Crefname{originallemma}{Original lemma}{Original lemmas}
\makeatletter
\newlabel{schemes-lemma-specialize-points}{{01J7}{source}{}{sourcepointer-5}{}}
\newlabel{schemes-lemma-specialize-points@cref}{{[originallemma][0][]01J7}{[1][0][]source}{}{}{}}
\makeatother
\crefname{originallemma}{original lemma}{original lemmas}
\Crefname{originallemma}{Original lemma}{Original lemmas}
\makeatletter
\newlabel{topology-lemma-dimension-function-catenary}{{02IA}{source}{}{sourcepointer-6}{}}
\newlabel{topology-lemma-dimension-function-catenary@cref}{{[originallemma][0][]02IA}{[1][0][]source}{}{}{}}
\makeatother
\title{A dimension function on a universally catenary locally Noetherian base pulls back to a dimension function on any locally finite type scheme over it}
\author{The Stacks Project Authors}
\date{stacks-project:morphisms.tex:local-snapshot-20261012; GNU FDL 1.2 or later}
\begin{document}
\maketitle
\noindent\textit{Editorial note.} All delivered mathematical text is the author's own, copied from the pinned Stacks Project chapter file pinned-original-source.tex; no statement or proof step was rewritten, reordered or removed.\par\medskip
\noindent\textit{Editorial note.} Following the multi-claim convention for this source, source\_url is the source-level URL https://stacks.math.columbia.edu/ and the source label lemma-dimension-function-propagates is carried in claim\_key and here rather than in a per-claim tag URL.\par\medskip
\noindent\textit{Editorial note.} The preamble is the chapter's shared Stacks Project preamble, retained as sources/stacks-morphisms/retained-preamble.tex: a line-for-line excerpt of the shared Stacks preamble staged with this snapshot. No macro was invented or redefined, and this card adds no preamble text of its own.\par\medskip
\noindent\textit{Editorial note.} Statement span 13656-13668 and proof span 13670-13738 come from the halfA pool card operations/pool-cards/stacks-morphisms-halfA.json and were re-read in the pinned chapter file before this card was written.\par\medskip
\noindent\textit{Editorial note.} The proof contains a three-line display computing the difference of the two values of the candidate dimension function; the display and the author's reference to the dimension formula are preserved unchanged.\par\medskip
\noindent\textit{Editorial note.} The author's notation for the candidate dimension function as both a function and a family indexed by the morphism is kept as written.\par\medskip
\noindent\textit{Editorial note.} 6 cross-references occur inside the counted spans; they are listed in external\_references as navigation pointers only.\par\medskip
\noindent\textit{Allowed standard prerequisites (editorial).}\begin{itemize}
\item The same-chapter criterion for universal catenarity to be local, and the characterization of locally finite type morphisms by finite type ring maps on affine charts. Cited at proof lines 13677-13681. Makes the coordinate ring of an affine chart universally catenary and the corresponding ring map of finite type.
\item The dimension formula for a finite type map of domains, cited from the algebra chapter as an equation, together with the description of points specializing to a given point through the spectrum of the local ring. Cited at proof lines 13719-13722 and 13689-13692. Supplies the identity relating transcendence degrees and dimensions of local rings used in the computation, and the reading of immediate specializations via the local ring's dimension.
\item For a dimension function on a catenary scheme the value on a generic point minus the value on a specialization equals the codimension of the closure, and that codimension is the dimension of the local ring at the specialization. Cited at proof lines 13724-13728. Turns the difference of the values of the dimension function on the base into the dimension of the local ring of the base at the point.
\end{itemize}\medskip
\ifdefined\hypertarget\hypertarget{sourcepointer-1}{}\fi
\noindent\textit{Source reference.} Original equation 02JV: focused-seeds/AGP/sources/stacks-snapshot-20261011/tags.txt, source lines 3350--3350. Same-chapter reference used in the retained author text at source line(s) 13725. It is supplied only as a navigation pointer so that the retained text compiles; the counted item delivers the author's own statement and the author's complete proof, and no part of the referenced item is copied in or substituted for the counted argument.\par\medskip
\ifdefined\hypertarget\hypertarget{sourcepointer-2}{}\fi
\noindent\textit{Source reference.} Original lemma 01T2: focused-seeds/AGP/sources/stacks-snapshot-20261011/tags.txt, source lines 2359--2359. Same-chapter reference used in the retained author text at source line(s) 2717, 2739, 2801, 2863, 3019, 3035, 3157, 3348, 3613, 3707, 3868, 5197, 5282, 6702, 9325, 9390, 9440, 11421, 13422, 13686. It is supplied only as a navigation pointer so that the retained text compiles; the counted item delivers the author's own statement and the author's complete proof, and no part of the referenced item is copied in or substituted for the counted argument.\par\medskip
\ifdefined\hypertarget\hypertarget{sourcepointer-3}{}\fi
\noindent\textit{Source reference.} Original lemma 02J9: focused-seeds/AGP/sources/stacks-snapshot-20261011/tags.txt, source lines 3328--3328. Same-chapter reference used in the retained author text at source line(s) 3185, 3189, 3223, 3452, 3457, 3484, 13684. It is supplied only as a navigation pointer so that the retained text compiles; the counted item delivers the author's own statement and the author's complete proof, and no part of the referenced item is copied in or substituted for the counted argument.\par\medskip
\ifdefined\hypertarget\hypertarget{sourcepointer-4}{}\fi
\noindent\textit{Source reference.} Original lemma 02IZ: focused-seeds/AGP/sources/stacks-snapshot-20261011/tags.txt, source lines 3318--3318. Same-chapter reference used in the retained author text at source line(s) 13732. It is supplied only as a navigation pointer so that the retained text compiles; the counted item delivers the author's own statement and the author's complete proof, and no part of the referenced item is copied in or substituted for the counted argument.\par\medskip
\ifdefined\hypertarget\hypertarget{sourcepointer-5}{}\fi
\noindent\textit{Source reference.} Original lemma 01J7: focused-seeds/AGP/sources/stacks-snapshot-20261011/tags.txt, source lines 2004--2004. Same-chapter reference used in the retained author text at source line(s) 1446, 12802, 13631, 13705, 13767. It is supplied only as a navigation pointer so that the retained text compiles; the counted item delivers the author's own statement and the author's complete proof, and no part of the referenced item is copied in or substituted for the counted argument.\par\medskip
\ifdefined\hypertarget\hypertarget{sourcepointer-6}{}\fi
\noindent\textit{Source reference.} Original lemma 02IA: focused-seeds/AGP/sources/stacks-snapshot-20261011/tags.txt, source lines 3293--3293. Same-chapter reference used in the retained author text at source line(s) 13729. It is supplied only as a navigation pointer so that the retained text compiles; the counted item delivers the author's own statement and the author's complete proof, and no part of the referenced item is copied in or substituted for the counted argument.\par\medskip
\begin{lemma}
\label{lemma-dimension-function-propagates}
Let $S$ be a locally Noetherian and universally catenary scheme.
Let $\delta : S \to \mathbf{Z}$ be a dimension function.
Let $f : X \to S$ be a morphism of schemes.
Assume $f$ locally of finite type.
Then the map
\begin{align*}
\delta = \delta_{X/S} : X & \longrightarrow \mathbf{Z} \\
x & \longmapsto \delta(f(x)) + \text{trdeg}_{\kappa(f(x))} \kappa(x)
\end{align*}
is a dimension function on $X$.
\end{lemma}

\begin{proof}
Let $f : X \to S$ be locally of finite type.
Let $x \leadsto y$, $x \not = y$ be a specialization in $X$.
We have to show that $\delta_{X/S}(x) > \delta_{X/S}(y)$ and
that $\delta_{X/S}(x) = \delta_{X/S}(y) + 1$ if $y$ is an
immediate specialization of $x$.

\medskip\noindent
Choose an affine open $V \subset S$ containing the image of
$y$ and choose an affine open $U \subset X$ mapping into $V$
and containing $y$. We may clearly replace $X$ by $U$ and
$S$ by $V$. Thus we may assume that $X = \Spec(A)$
and $S = \Spec(R)$ and that $f$ is given by a ring
map $R \to A$. The ring $R$ is universally catenary
(Lemma \ref{lemma-universally-catenary-local})
and the map $R \to A$ is of finite type
(Lemma \ref{lemma-locally-finite-type-characterize}).

\medskip\noindent
Let $\mathfrak q \subset A$ be the prime ideal corresponding
to the point $x$ and let $\mathfrak p \subset R$ be the prime
ideal corresponding to $f(x)$. The restriction $\delta'$ of $\delta$
to $S' = \Spec(R/\mathfrak p) \subset S$ is a dimension
function. The ring $R/\mathfrak p$ is universally catenary.
The restriction of $\delta_{X/S}$ to $X' = \Spec(A/\mathfrak q)$
is clearly equal to the function $\delta_{X'/S'}$ constructed
using the dimension function $\delta'$. Hence we may assume
in addition to the above that $R \subset A$ are domains, in
other words that $X$ and $S$ are integral schemes, and that
$x$ is the generic point of $X$ and $f(x)$ is the generic point of $S$.

\medskip\noindent
Note that $\mathcal{O}_{X, x} = R(X)$ and that since $x \leadsto y$,
$x \not = y$, the spectrum of $\mathcal{O}_{X, y}$ has at least
two points
(Schemes, Lemma \ref{schemes-lemma-specialize-points})
hence $\dim(\mathcal{O}_{X, y}) > 0$ .
If $y$ is an immediate specialization of $x$, then
$\Spec(\mathcal{O}_{X, y}) = \{x, y\}$ and
$\dim(\mathcal{O}_{X, y}) = 1$.

\medskip\noindent
Write $s = f(x)$ and $t = f(y)$. We compute
\begin{align*}
\delta_{X/S}(x) - \delta_{X/S}(y)
& =
\delta(s) + \text{trdeg}_{\kappa(s)} \kappa(x)
- \delta(t) - \text{trdeg}_{\kappa(t)} \kappa(y) \\
& =
\delta(s) - \delta(t) +
\text{trdeg}_{R(S)} R(X) - \text{trdeg}_{\kappa(t)} \kappa(y) \\
& =
\delta(s) - \delta(t) + \dim(\mathcal{O}_{X, y})
- \dim(\mathcal{O}_{S, t})
\end{align*}
where we use equality in (\ref{equation-dimension-formula}) in
the last step. Since $\delta$ is a dimension function on the scheme $S$
and $s \in S$ is the generic point, the difference
$\delta(s) - \delta(t)$ is equal to $\text{codim}(\overline{\{t\}}, S)$ by
Topology, Lemma \ref{topology-lemma-dimension-function-catenary}.
This is equal to $\dim(\mathcal{O}_{S, t})$ by
Properties, Lemma
\ref{properties-lemma-codimension-local-ring}.
Hence we conclude that
$$
\delta_{X/S}(x) - \delta_{X/S}(y) = \dim(\mathcal{O}_{X, y})
$$
and the lemma follows from what we said above about $\dim(\mathcal{O}_{X, y})$.
\end{proof}

\end{document}
